% !TeX root = ../PIFB2.tex
\section{Consciousness: An Interpretive Stance}\label{sec:consciousness}

The theory developed here concerns statistical agents, communication and retained laws. It does not attach experience to a probability distribution by mathematical necessity. This section records a possible cognitive interpretation and the additional empirical question it would create. No theorem in the monograph depends on that interpretation.

A cognitive frame can represent how a system presents and communicates its variables. Its invariant content lies in the laws, records, interfaces and responses that survive a passive change of presentation. A relabeling of internal coordinates changes none of that content. It therefore cannot, by itself, model an altered experience, a perceptual intervention or a change of cognitive architecture. Such a proposal must change some operational part of the model, for example the record-access channel, the represented latent variables, a prior or the learned interface. It must then explain how that change would be detected.

The covariant Fisher tensor is a candidate descriptor of statistical sensitivity, and the collective law can describe integration of information unavailable in individual marginals. Neither object supplies a correspondence between a geometric value and a particular phenomenal quality. Two systems can share a Fisher tensor while differing in higher statistical structure, history and observation mechanisms. Conversely, a coordinate change can alter the matrix entries of a representation without altering any experience-related observable the model actually defines. These distinctions rule out a formal solution to the inverted-spectrum problem by coordinate relabeling alone.

Meta-agent formation supplies no criterion of consciousness either. A retained collective may be coherent, persistent and predictively closed, with or without explicitly retained memory. Those are properties of its law and dynamics. To propose them as necessary or sufficient for conscious experience would require an independent operational criterion and cases that distinguish the proposal from rival explanations. Recursion or cross-scale feedback does not establish that criterion merely because it occurs at several levels.

The participatory interpretation is consequently compatible with several positions about the relation between inference and experience. It does not derive panpsychism, eliminate an observer-independent reality, or settle why information processing should be accompanied by experience. The narrower research opportunity is to construct competing models of report, discrimination and access, and test which collective variables improve their predictions under intervention. Such results could constrain a cognitive interpretation while leaving the broader philosophical question open.
